\documentclass[12pt,a4paper]{article}

\usepackage[utf8]{inputenc}
\usepackage{amsmath}
\usepackage{graphicx}
\usepackage{float}
\usepackage{hyperref}
\usepackage{geometry}
\geometry{margin=1in}

\begin{document}

\begin{titlepage}
    \centering
    \vspace*{6.5cm}
    {\LARGE \textbf{Network Intrusion Detection System} \\[0.5em]
        \textbf{Powered by Machine Learning}}\\[2cm]
    {\large Phase 3: Data Preparation}\\[2cm]
    {\large Hajer JEMAA \& Amira BOUAFIF}\\[0.5cm]
    \vfill
    {\large 2025-2026}\\[1cm]
\end{titlepage}

\section{Introduction}

This report presents the Data Preparation phase of the \textbf{CRISP-DM}methodology applied to the \textbf{CIC-IDS2017} dataset. The objective of this phase is to transform the \textbf{CIC-IDS2017} flow-level data into a clean, consistent, and model-ready dataset. This encompasses two main stages: data cleaning and feature selection, which address encoding inconsistencies, infinite and missing values, duplicate rows, and redundant features; followed by data preprocessing, which covers label encoding, stratified K-Fold splitting, class imbalance handling, and feature normalization.

\section{Methodology}
This report follows the \textbf{CRISP-DM} (Cross-Industry Standard Process for Data Mining) methodology, a widely used framework for structuring machine learning projects. The Data Preparation phase involves cleaning and transforming the data to make it suitable for modeling.

\section{Column and Label Normalisation}
This section addresses the issues of inconsistent column naming and label encoding found in the Data Understanding phase.

\subsection{Whitespace Correction}
Although the whitespace correction is applied during the Data Understanding phase, it is not persisted. It is only applied to the in-memory DataFrame and not to the source files. As is noted in that report, this issue is an artefact of CICFlowMeter's CSV export format. Reapplying it here ensures that the correction is included in the cleaned dataset saved at the end of this phase, so that all subsequent steps work on already-corrected data.

\subsection{Label Encoding Correction}
As is noted in the Data Understanding report, three label values in the Thursday-AM file contain encoding issues. As with the whitespace correction, the fix is reapplied here to ensure it is included in the cleaned dataset saved at the end of this phase.

\subsection{Concatenation}
After applying the whitespace and encoding corrections, the 8 CSV files are concatenated into a single DataFrame \texttt{df\_all} of \textbf{2,830,743 rows and 80 columns} (78 numeric features, the target label, and the \texttt{source\_file} column, noting that one duplicate column is still present at this stage). A \texttt{source\_file} column is added to track the origin of,each row, which is used later to audit fold composition and investigate per-session feature drift identified in the Data Understanding phase. This ensures that all subsequent operations, including cleaning, duplicate detection, and splitting, are applied uniformly across the entire dataset, so that no cross-file duplicates are missed and the split reflects the true class distribution.

\section{Inf/NaN Cleaning}
In the Data Understanding phase, 2,867 rows, representing approximately 0.1\% of
the dataset, are found to contain infinite or missing values. The cause for those infinite and missing values is explained in the Data Understanding report. These rows are dropped rather than imputed, as the subsequent analysis reveals that rare classes such as Heartbleed, Web Attack - Sql Injection, and Infiltration are not affected, and the overall impact on the dataset is minimal. Table~\ref{tab:nan_inf} presents the distribution of NaN and Inf values across classes, showing that the majority of affected rows belong to the BENIGN class, which is well-represented in the dataset and can tolerate some loss without jeopardising model performance.

\begin{table}[H]
    \centering
    \caption{Class level Impact of Inf/NaN values}
    \vskip0.1cm
    \label{tab:nan_inf}
    \begin{tabular}{|l|r|r|r|r|}
        \hline
        \textbf{Class} & \textbf{Total Rows} & \textbf{\% of The Dataset} & \textbf{Rows Affected} & \textbf{\% of Affected Rows} \\
        \hline
        BENIGN         & 2,273,097           & 80.300\%                   & 1,777                  & 61.981\%                     \\
        \hline
        DoS Hulk       & 231,073             & 8.163\%                    & 949                    & 33.101\%                     \\
        \hline
        PortScan       & 158,930             & 5.614\%                    & 126                    & 4.395\%                      \\
        \hline
        Bot            & 1,966               & 0.069\%                    & 10                     & 0.349\%                      \\
        \hline
        FTP-Patator    & 7,938               & 0.280\%                    & 3                      & 0.105\%                      \\
        \hline
        DDoS           & 128,027             & 4.523\%                    & 2                      & 0.070\%                      \\
        \hline
    \end{tabular}
\end{table}

After dropping the rows containing Inf/NaN values, the dataset is reduced to
2,827,876 rows, losing only 0.1013\% of the data.

\section{Duplicate Removal}
This section addresses the issue of duplicate rows identified in the Data Understanding phase. Two types of duplicates are identified: contradictory duplicates (same features, different labels) and exact duplicates (same features and labels).

\subsection{Contradictory Duplicates}
In the Data Understanding phase, a total of 698 contradictory patterns are found affecting 7,020 rows. After reapplying the check following the Inf/NaN cleaning, one contradictory pattern is resolved, leaving 697 patterns with a total of 6,666 affected rows. Table~\ref{tab:con_pat} shows a section of the largest contradictory patterns.

\begin{table}[H]
    \centering
    \caption{Contradictory duplicate patterns samples}
    \vskip0.1cm
    \label{tab:con_pat}
    \begin{tabular}{|l|r|l|l|}
        \hline
        \textbf{No.} & \textbf{Pattern size} & \textbf{Class 1} & \textbf{Class 2}  \\
        \hline
        1       &3362  & BENIGN (2)            & DoS Hulk (3360)  \\
        \hline
        2       &164  & BENIGN (1)            & DoS Hulk (163)   \\
        \hline
        3       &133  & BENIGN (2)            & DoS Hulk (131)   \\
        \hline
        4       &81  & BENIGN (1)            & DoS Hulk (80)    \\
        \hline
        5       &61  & BENIGN (1)            & DoS Hulk (60)    \\
        \hline
        6       &54  & BENIGN (1)            & DoS Hulk (53)    \\
        \hline
        7       &48  & BENIGN (1)            & DoS Hulk (47)    \\
        \hline
    \end{tabular}
\end{table}

As shown in Table~\ref{tab:con_pat}, each contradictory pattern involves a conflict between two classes: one BENIGN and one attack class. As noted in the Data Understanding report, the majority of affected rows belong to the DoS Hulk class. As the true label cannot be determined from the feature vectors alone and no rare class is affected, all 6,666 rows are dropped to prevent introducing noise and confusion into the model training process. This results in a reduction of approximately 0.24\% of the dataset, leaving 2,821,210 rows.

\subsection{Exact Duplicates}
In the Data Understanding phase, a total of 308,381 exact duplicate rows are identified. After reapplying the check following the Inf/NaN cleaning and contradictory duplicate removal, it went down by 6,575 rows to 301,806 making approximately 10.7\% of the dataset. As the exact duplicates do not introduce noise but only redundancy and no rare class is affected, they are removed by keeping one instance of each duplicate pattern leaving 2,519,404 rows in the dataset.       
\section{Feature Removal}
This section addresses the issue of redundant features identified in the Data Understanding phase. Three types are identified: duplicate columns, constant features, and near-constant features.
\subsection{Duplicate Column}
In the Data Understanding phase, one duplicate column is identified: \texttt{Fwd Header Length.1} is an exact duplicate of \texttt{Fwd Header Length}. The duplicate column is dropped, leaving 79 columns (78 features, the target label, and the \texttt{source\_file} column) in the dataset.
\subsection{Constant Features}
In the Data Understanding phase, 8 constant columns are identified, comprising 6 bulk transfer features and two flag count features, all with a unique value of 0. As these columns carry no discriminative information, they are dropped, leaving 71 columns.
\subsection{Near-Constant Features}
In the Data Understanding phase, 4 near-constant features are identified, each with only two unique values and a dominant value of 0 with a dominance exceeding 99.97\%. As with the constants features, these features carry no discriminative information for classification. Furthermore, it is confirmed that no attack class, including the rare classes, has any non-zero values for these features. They are therefore dropped, leaving 67 columns remaining. 
\section{Correlation-Based Feature Selection}
This section 
\subsection{Recomputed Correlation Analysis}

\subsection{Mutual Information Scoring}

\subsection{Decision on Correlated Pairs}

\subsection{Domain Knowledge Overrides}

\subsection{Final Feature Set}

\section{Label Encoding}

The target column \texttt{Label} contains string-valued class names such as
\texttt{BENIGN}, \texttt{DDoS}, and \texttt{PortScan}. XGBoost, the chosen model for this project, requires numeric targets. Label encoding
is therefore a mandatory preprocessing step before any modelling or evaluation
is performed.

\subsection{Encoding Scheme}

The \texttt{LabelEncoder} from \texttt{sklearn.preprocessing} is applied to the
\texttt{Label} column. It assigns one integer identifier per class in lexicographic
order, producing a contiguous mapping from zero to fourteen across the fifteen classes
present in the dataset. Table~\ref{tab:label_encoding} presents the resulting
class-to-identifier mapping.

\begin{table}[H]
    \centering
    \caption{Class-to-identifier mapping produced by the label encoder.}
    \label{tab:label_encoding}
    \begin{tabular}{|r|l|}
        \hline
        \textbf{Identifier} & \textbf{Class name}        \\
        \hline
        0                   & BENIGN                     \\
        \hline
        1                   & Bot                        \\
        \hline
        2                   & DDoS                       \\
        \hline
        3                   & DoS GoldenEye              \\
        \hline
        4                   & DoS Hulk                   \\
        \hline
        5                   & DoS Slowhttptest           \\
        \hline
        6                   & DoS slowloris              \\
        \hline
        7                   & FTP-Patator                \\
        \hline
        8                   & Heartbleed                 \\
        \hline
        9                   & Infiltration               \\
        \hline
        10                  & PortScan                   \\
        \hline
        11                  & SSH-Patator                \\
        \hline
        12                  & Web Attack - Brute Force   \\
        \hline
        13                  & Web Attack - Sql Injection \\
        \hline
        14                  & Web Attack - XSS           \\
        \hline
    \end{tabular}
\end{table}

As shown in Table~\ref{tab:label_encoding}, the assignment of identifiers follows
strict lexicographic order. It is essential to note that the resulting integer codes
carry no ordinal meaning whatsoever: identifier 10 does not denote a class that is
more severe or more frequent than identifier 2. The codes are pure nominal identifiers
whose numeric values are an artefact of the alphabetical ordering of class names.
Any interpretation that assigns magnitude or rank to these values would be incorrect.

\subsection{Design Decisions}

Two design decisions govern the application of the encoder and warrant explicit
justification.

The encoder is fitted once on the full combined dataset, prior to any train/test
splitting. This guarantees a single, consistent class-to-identifier mapping across
all downstream folds. If the encoder were fitted independently within each training
fold, the lexicographic mapping would remain stable in this case because all fifteen
classes are present in the full dataset, however fitting within folds would
introduce a dependency on the fold composition and would risk identifier drift if any fold
were to lack a class, which is a realistic scenario for the three Critical classes
(Heartbleed with 11 samples, Infiltration with 36, and Web Attack - Sql Injection
with 21) even under stratified sampling. Fitting on the full dataset before splitting
eliminates this risk and ensures that training and evaluation always operate on
identical identifier spaces.

The encoded values are stored in a new column \texttt{label\_multi}, added alongside
the original \texttt{Label} column rather than replacing it. The original string
column is preserved so that the dataset remains human-readable throughout the
preparation pipeline and so that encoded predictions can be validated against
ground-truth class names at any stage without requiring a reverse lookup operation.

\subsection{Mapping Persistence}

The class-to-identifier mapping is serialised to a JSON file at
\texttt{data/prepared/label\_mapping.json}, associating each integer identifier
with its corresponding class name. This file serves two purposes. First, it ensures
reproducibility: any future process that produces or consumes encoded labels can
reload the mapping and verify that identifiers are consistent with those used during
training, without requiring access to the fitted encoder object. Second, it enables
interpretability during evaluation, as predicted integer identifiers can be decoded
back to human-readable class names for reporting and error analysis.

\section{Conclusion}

\end{document}